\documentclass[
   ngerman          % neue deutsche Rechtschreibung
  ,a4paper          % Papiergrösse
  ,12pt             % Schriftgrösse
  ,pdftex
%  ,disable         % Todo-Markierungen auschalten
]{scrreprt}

\usepackage[utf8]{inputenc}        % UTF-8 codierte Dateien
                                   % Dieses Dokument ist unter Unix erstellt, daher
                                   % wird diese Input-Codierung benutzt.
\usepackage[T1]{fontenc}

\usepackage{bericht}

\usepackage{gradientframe}

\usepackage{colortbl}

\usepackage{rotating}

\usepackage{eurosym}

\usepackage{paralist}

\usepackage{enumitem}

\usepackage{rotating}

\usepackage{multirow}

\usepackage{multicol} 

\usepackage{arydshln}

\usepackage{color}

\usepackage[onehalfspacing]{setspace}

\usepackage{listings}             % Include the listings-package for Code
%%\lstset{language=[Sharp]C} % Set your language (you can change the language for each code-block optionally)

\lstset{
language=csh,
basicstyle=\normalsize\ttfamily,
numbers=left,
numberstyle=\tiny,
numbersep=5pt,
numberstyle=\scriptsize,
tabsize=2,
extendedchars=true,
breaklines=true,
frame=b,
stringstyle=\color{blue}\ttfamily,
showspaces=false,
showtabs=false,
xleftmargin=17pt,
framexleftmargin=17pt,
framexrightmargin=5pt,
framexbottommargin=4pt,
commentstyle=\color{green},
morecomment=[l]{//}, %use comment-line-style!
morecomment=[s]{/*}{*/}, %for multiline comments
showstringspaces=false,
morekeywords={ abstract, event, new, struct,
as, explicit, null, switch,
base, extern, object, this,
bool, false, operator, throw,
break, finally, out, true,
byte, fixed, override, try,
case, float, params, typeof,
catch, for, private, uint,
char, foreach, protected, ulong,
checked, goto, public, unchecked,
class, if, readonly, unsafe,
const, implicit, ref, ushort,
continue, in, return, using,
decimal, int, sbyte, virtual,
default, interface, sealed, volatile,
delegate, internal, short, void,
do, is, sizeof, while,
double, lock, stackalloc,
else, long, static,
enum, namespace, string},
keywordstyle=\color{red},
identifierstyle=\color{black},
backgroundcolor=\color{white},
}

\definecolor{editorGray}{rgb}{0.95, 0.95, 0.95}
\definecolor{editorOcher}{rgb}{1, 0.5, 0} % #FF7F00 -> rgb(239, 169, 0)
\definecolor{editorGreen}{rgb}{0, 0.5, 0} % #007C00 -> rgb(0, 124, 0)

\usepackage{minted}

\usepackage{float}
\restylefloat{figure}


\definecolor{javared}{rgb}{0.6,0,0} % for strings
\definecolor{javagreen}{rgb}{0.25,0.5,0.35} % comments
\definecolor{javapurple}{rgb}{0.5,0,0.35} % keywords
\definecolor{javadocblue}{rgb}{0.25,0.35,0.75} % javadoc

\renewcaptionname{ngerman}{\abstractname}{Kurzzusammenfassung}


%%%%%%%%%%%%%%%%%%%%%%%%%%%%%%%%%%%%%%%%%%%%%%%%%%%%%%%%%%%%%%%%%%%%%%%%%%%%%%%
%% Angaben zur Arbeit
%%%%%%%%%%%%%%%%%%%%%%%%%%%%%%%%%%%%%%%%%%%%%%%%%%%%%%%%%%%%%%%%%%%%%%%%%%%%%%%

\newcommand{\Autor}{Avelina Ott}
\newcommand{\MatrikelNummer}{4386414}
\newcommand{\Kursbezeichnung}{TINF15B4}

\newcommand{\FirmenName}{WebCommerce GmbH}
\newcommand{\FirmenStadt}{Offenburg}
\newcommand{\FirmenLogoDeckblatt}{\includegraphics[width=3cm]{bilder/DHBW.png}}
\newcommand{\BetreuerFirma}{Markus Schill}
\newcommand{\BetreuerDHBW}{Dr. Werner Geiger}



%%%%%%%%%%%%%%%%%%%%%%%%%%%%%%%%%%%%%%%%%%%%%%%%%%%%%%%%%%%%%%%%%%%%%%%%%%%%%%
%%%% Befehle in der Arbeit
%%%%%%%%%%%%%%%%%%%%%%%%%%%%%%%%%%%%%%%%%%%%%%%%%%%%%%%%%%%%%%%%%%%%%%%%%%%%%%
\newcommand{\edith}{EDITH}
\newcommand{\cms}{Content-Management-System}
\newcommand{\cmsEdith}{CMS EDITH}

\newcommand{\exkurs}[2]{
\begin{leftbar}
	\textbf{Exkurs - } \textit{#1}: #2
\end{leftbar}
}


%%%%%%%%%%%%%%%%%%%%%%%%%%%%%%%%%%%%%%%%%%%%%%%%%%%%%%%%%%%%%%%%%%%%%%%%%%%%%%%%%%%%%

\newcommand{\Was}{Bachelorarbeit}
% Wird auf dem Deckblatt in der Erklärung benutzt

%%%%%%%%%%%%%%%%%%%%%%%%%%%%%%%%%%%%%%%%%%%%%%%%%%%%%%%%%%%%%%%%%%%%%%%%%%%%%%%%%%%%%




% Benutzt man das "biblatex"-Paket, dann muß das hier stehen:
% siehe auch die mit BIBLATEX markierten Zeilen in bericht.sty
\bibliography{literatur}

\begin{document}
\pagenumbering{Roman}
%%%%%%%%%%%%%%%%%%%%%%%%%%%%%%%%%%%%%%%%%%%%%%%%%%%%%%%%%%%%%%%%%%%%%%%%%%%%%%%

\newcommand{\Titel}{Web-to-Print: Export von Terminen in eine kundenspezifische DOCX-Datei}
\newcommand{\AbgabeDatum}{27. August 2018}

\newcommand{\Dauer}{12 Wochen}

\newcommand{\Abschluss}{Bachelor of Science}

\newcommand{\Studiengang}{Angewandte Informatik}

\hypersetup{%%
  pdfauthor={\Autor},
  pdftitle={\Titel},
  pdfsubject={\Was}
}





\begin{titlepage}
\begin{center}
\vspace*{-2cm}
%\FirmenLogoDeckblatt
\hfill\includegraphics[width=4cm]{bilder/DHBW}\\[2cm]
{\Huge \Titel}\\[2cm]
{\Huge\scshape \Was}\\[2cm]
{\large für die Prüfung zum}\\[0.5cm]
{\Large \Abschluss}\\[0.5cm]
{\large des Studienganges \Studiengang}\\[0.5cm]
{\large an der}\\[0.5cm]
{\large Dualen Hochschule Baden-Württemberg Karlsruhe}\\[0.5cm]
{\large von}\\[0.5cm]
{\large\bfseries \Autor}\\[1cm]
{\large \AbgabeDatum}
\vfill
\end{center}
\begin{tabular}{l@{\hspace{2cm}}l}
Bearbeitungszeitraum	         & \Dauer 			\\
%Matrikelnummer	                 & \MatrikelNummer	\\
Kurs			         & \Kursbezeichnung		\\
Ausbildungsfirma	         & \FirmenName			\\
			         %& \FirmenStadt			\\
Betreuer der Ausbildungsfirma	 & \BetreuerFirma		\\
Gutachter/in der Studienakademie	 & \BetreuerDHBW		\\
\end{tabular}
\end{titlepage}

%%%%%%%%%%%%%%%%%%%%%%%%%%%%%%%%%%%%%%%%%%%%%%%%%%%%%%%%%%%%%%%%%%%%%%%%%%%%%%%

% Nur für Bachelorarbeiten einfügen:

\selectlanguage{ngerman}
%\begin{abstract}
%\end{abstract}


\newpage
\addchap{Eidesstattliche Erklärung}
Gemäß § 5(3) der „Studien- und Prüfungsordnung DHBW Technik“ vom 29. 9. 2015 versichere ich hiermit, dass ich meine Bachelorarbeit mit dem Thema:
   \begin{quote}
     „\textit{\Titel}“
    \end{quote}
selbstständig verfasst und keine anderen als die angegebenen Quellen und Hilfsmittel benutzt habe. Ich versichere zudem, dass die eingereichte elektronische Fassung mit der gedruckten Fassung übereinstimmt.


\vspace{3cm}
\noindent
\underline{\hspace{10cm}}\\\\
Karlsruhe, 27.08.2018\hspace{4cm}\\\\\\

\tableofcontents           % Inhaltsverzeichnis hier ausgeben
\listoffigures             % Liste der Abbildungen
\listoftables              % Liste der Tabellen
\lstlistoflistings         % Liste der Listings


\lstdefinelanguage{Velocity}
{
  %basicstyle=\ttfamily,
  morestring=[s]{"}{"},
  morekeywords={set, foreach, end, if, else},
  otherkeywords={\#, ., get},
  %moredelim=[s][\color{black}]{>}{<},
  moredelim=[s][\color{red}]{\$}{\ },
  %morecomment=[s][\color{red}]{.}{.},
  %moredelim=[s][\color{red}]{\$}{)},
  %moredelim=[s][\color{red}]{\$}{.},
  morecomment=[s][\color{green}]{.}{(},
  commentstyle=\color{green},
  stringstyle=\color{blue},
  keywordstyle=\color{green}
}

\lstdefinelanguage{HTML5}{
        language=html,
        sensitive=true, 
        alsoletter={<>=-},
        otherkeywords={
        % HTML tags
        <html>, <head>, <title>, </title>, <meta, />, </head>, <body>,
        <canvas, \/canvas>, <script>, </script>, <td>, </td>, <tr>, </tr>, </body>, </html>, <!, html>, <style>, </style>, ><
        },  
        ndkeywords={
        % General
        =,
        % HTML attributes
        charset=, id=, width=, height=,
        % CSS properties
        border:, transform:, -moz-transform:, transition-duration:, transition-property:, transition-timing-function:
        },  
        morecomment=[s]{<!--}{-->},
        tag=[s]
}





%\lstset{language=Java,
%	%basicstyle=\fontsize{10}{13}\ttfamily,
%	keywordstyle=\color{javapurple}\bfseries,
%	stringstyle=\color{javared},
%	commentstyle=\color{javagreen},
%	morecomment=[s][\color{javadocblue}]{/**}{*/},
%	numbers=left,
%	numberstyle=\tiny\color{black},
%	stepnumber=1,
%	numbersep=5pt,
%	tabsize=2,
%	showspaces=false,
%	showstringspaces=false,
%	breaklines=true,
%	breakautoindent=true,
%	morekeywords={set, foreach, end}
%}


\lstset{
numbers = left,
frame = single,
framexleftmargin=15pt,
breaklines=true
}

% Jetzt kommt der "eigentliche" Text
\include{abk}              % Abkürzungsverzeichnis

\pagenumbering{arabic}

\chapter{Einleitung}

\section{Problemstellung} 


\section{Motivation} 


\section{Zielsetzung} 



\section{Vorgehensweise} 


\chapter{State of the Art}
--> welche wissenschaftlichen Kenntnisse gibt es bereits, bei der Produktentwicklung berücksichtigt werden sollten?
    --> mögliche Lösungen schon bereits enthalten -> wie wird es bereits woanders umgesetzt?

--> Welche Methode wird angewandt?
    --> hier muss so was rein wie verschiedene Arten von Webt-To-Print, worin sie sich entscheiden, welche Definition verwendet wird etc.

--> hier eventuell ein Vergleich mit bereits existierenden Lösungen von anderen CMS

\chapter{Anforderungsanalyse} \label{anforderungsanalyse}

--> Welche Spezifikation gilt für das Produkt und wie wurde die Spezifikation ermittelt?

--> Beschreibung des Produktes anhand der wichtigsten Funktionen und Leistungen? --> Eventuell auch Gebrauchsnutzen?

\section{Funktionale Anforderungen} 



\section{Nicht funktionale Anforderungen}






\chapter{Planung und Konzept} \label{planung}



\section{Projektplan} %Michael



\section{Vorgehensweise} %Michael

\chapter{Implementierung}\label{implem} %Niko

\chapter{Evaluierung} \label{evaluierung}

\subsection{Metriken}\label{metriken}
--> wie wird das Produkt erprobt? 
--> Welches Ergebnis zeigt sich?
--> Anhand welcher Metriken wird evaluiert?

\subsection{Ergebnis} \label{ergebnis}
--> Welche Konsequenzen ergeben sich aus der Erprobung und Evaluierung des Produktes?
\chapter{Zusammenfassung} \label{zusammenfassung}

\section{Fazit} \label{fazit}

\section{Ausblick}\label{Ausblick}


%\include{anhang}
\pagenumbering{Roman}

\addcontentsline{toc}{chapter}{Literaturverzeichnis}
\printbibliography



% Ab hier beginnt der Anhang
\appendix
%\addcontentsline{toc}{chapter}{Anhang}
%\appendix
%\appendixtoc
%\newpage

\renewcommand*{\thesection}{\Alph{section}}
\chapter*{\appendixname}
%\include{anhang}
%\begin{enumerate}[label=\bfseries \LARGE \Alph*, ref=\Alph*]
%	\item Bilder \label{app:bilder}
%   \item Umfrage \label{app:umfrage}
%  \item manueller Testplan \label{app:maTestplan}
%\end{enumerate}




% Haben Sie das "biblatex"-Paket nicht installiert, benutzen Sie folgendes:
% Ohne das "biblatex"-Paket (s. bericht.sty) produziert folgendes
% "deutsche" Zitate in Literaturverzeichnissen gemaß der Norm DIN 1505,
% Teil 2 vom Jan. 1984.
% Die Zitatmarken werden alphabetisch nach Verfassern
% sortiert und sind durch abgekürzte Verfasserbuchstaben plus
% Erscheinungsjahr in eckigen Klammern gekennzeichnet.

%\bibliographystyle{alphadin}
%\bibliography{literatur}

%%%%%%%%%%%%%%%%%%%%%%%%%%%%%%%%%%%%%%%
% BIBLATEX
% Benutzt man das "biblatex"-Paket, muß man folgendes schreiben:
%\def\refname{Literaturverzeichnis}


%%%%%%%%%%%%%%%%%%%%%%%%%%%%%%%%%%%%%%%





\end{document}
